\documentclass[10pt,a4paper]{article}
\usepackage[T1]{fontenc}
\usepackage[utf8]{inputenc}
\usepackage[brazilian]{babel}
\usepackage[a4paper,margin=1.6cm]{geometry}
\usepackage{booktabs}
\usepackage{xcolor}
\usepackage{microtype}
\usepackage[hidelinks]{hyperref}
\pagestyle{empty}
\setlength{\parindent}{0pt}
\setlength{\parskip}{3pt}
\definecolor{chuva}{HTML}{1B7F9E}
\definecolor{cinza}{HTML}{666666}
\newcommand{\barra}[1]{\textcolor{chuva}{\rule{#1mm}{2.4pt}}}
\begin{document}

{\Large\bfseries Registro Unesp --- climatologia das médias}\\[1pt]
{\small Estação CIIAGRO/IAC 109 \,\textbullet\, 191 meses
medidos (2009-09 a 2026-10) \,\textbullet\, \textbf{14 anos com os 12 meses fechados}}

\vspace{3pt}
\colorbox{cinza!8}{\parbox{\dimexpr\textwidth-2\fboxsep}{\small
\textbf{Leia antes de usar:} os números abaixo são \textbf{médias de vários anos},
calculadas a partir das medições mensais do CIIAGRO. \textbf{Não é previsão do tempo}
para amanhã nem para este ano: é o padrão médio da estação, para apoiar decisões de
calendário agrícola, manejo de solo e planejamento da propriedade.}}

\vspace{5pt}
\begin{tabular}{@{}p{0.31\textwidth}p{0.31\textwidth}p{0.31\textwidth}@{}}
\textbf{Chuva anual média} & \textbf{Mês mais chuvoso} & \textbf{Mês mais seco} \\
{\LARGE 1.310} mm/ano &
{\LARGE 191} mm & {\LARGE 58} mm \\
{\small em todo o ano} & {\small em janeiro} &
{\small em julho} \\
\end{tabular}

\vspace{3pt}
\begin{tabular}{@{}p{0.31\textwidth}p{0.31\textwidth}p{0.31\textwidth}@{}}
\textbf{Temperatura média} & \textbf{Mês mais quente} & \textbf{Mês mais frio} \\
{\LARGE 21,3} °C & 26,1 °C &
16,6 °C \\
{\small em todo o ano} & {\small em janeiro} &
{\small em julho} \\
\end{tabular}

\vspace{6pt}
{\bfseries Chuva média por mês do ano}
\vspace{1pt}
{\small
\begin{tabular}{@{}l r l r r r@{}}
\toprule
mês & chuva (mm) & perfil & t mín & t méd & t máx \\
\midrule
janeiro    &   191 & \barra{26.0}   & 20,4 & 26,1 & 31,8 \\
fevereiro  &   175 & \barra{23.8}   & 19,9 & 25,5 & 27,9 \\
março      &   141 & \barra{19.2}   & 19,2 & 24,0 & 25,8 \\
abril      &    84 & \barra{11.5}   & 17,0 & 22,1 & 24,6 \\
maio       &    83 & \barra{11.3}   & 13,4 & 18,7 & 24,0 \\
junho      &   101 & \barra{13.8}   & 12,2 & 17,0 & 21,8 \\
julho      &    58 & \barra{7.9}    & 11,4 & 16,6 & 21,8 \\
agosto     &    64 & \barra{8.7}    & 12,4 & 17,5 & 22,7 \\
setembro   &    76 & \barra{10.4}   & 15,1 & 20,9 & 24,2 \\
outubro    &   123 & \barra{16.8}   & 15,2 & 21,4 & 25,0 \\
novembro   &   109 & \barra{14.8}   & 17,0 & 22,0 & 26,9 \\
dezembro   &   146 & \barra{19.9}   & 19,1 & 24,4 & 29,7 \\
\midrule
ano & 1.310 & & & 21,3 & \\
\bottomrule
\end{tabular}}

\vspace{5pt}
{\bfseries Chuva por estação do ano}
\vspace{1pt}
{\small
\begin{tabular}{@{}l l r r r l l@{}}
\toprule
estação & meses & total (mm) & \% do ano & média/mês & maior mês & menor mês \\
\midrule
Verão      & dez--fev &   511 &  39\% &  170 & 191 (jan) & 146 (dez) \\
Outono     & mar--mai &   308 &  23\% &  103 & 141 (mar) & 83 (mai) \\
Inverno    & jun--ago &   222 &  17\% &   74 & 101 (jun) & 58 (jul) \\
Primavera  & set--nov &   308 &  24\% &  103 & 123 (out) & 76 (set) \\
\bottomrule
\end{tabular}}

\vspace{3pt}
{\small\textbf{Antes dos solstícios} (acumulado desde janeiro)}

{\small
inverno: janeiro a maio somam 673 mm (51\% do ano) \quad\textbullet\quad
verão: junho a novembro somam 530 mm (41\% do ano)
\\ Dezembro fica depois do solstício de verão e entra no total anual, não nos
acumulados acima.}

\vspace{5pt}
{\bfseries Como usar no planejamento nesta estação}
\small\begin{itemize}\setlength{\itemsep}{1pt}
  \item \textbf{Amostragem de solo e preparo mecânico}: Em julho, agosto e setembro a chuva média fica entre 58 e 76 mm/mês — é a janela de menor encharcamento para retirar terra para laudo, arar e incorporar calcário com antecedência.
  \item \textbf{Calagem em cobertura (pasto e perene já formado)}: A partir de outubro, quando a média volta a passar de 120 mm/mês: aí há água para o corretivo descer no perfil sem enxurrada.
  \item \textbf{Plantio e transplante}: O verão (dezembro a fevereiro) concentra 511 mm, 39\% da chuva do ano — favorece o pegamento de mudas sem irrigação pesada; em áreas declivosas, reforce a conservação de solo nesse período.
  \item \textbf{Calor, frio e fitossanidade}: janeiro tem a maior média (26,1 °C) e julho a menor (16,6 °C, com mínima média de 11,4 °C). Antecipe o monitoramento de doenças fúngicas nos meses quentes e úmidos e proteja culturas sensíveis ao frio nos mais frios.
\end{itemize}

\vspace{2pt}
{\bfseries Vale para qualquer município da região}
\footnotesize\begin{itemize}\setlength{\itemsep}{1pt}
  \item \textbf{Média não é previsão}: Descrevem o padrão de vários anos, não o tempo deste ano: use para escolher a janela de cada operação, e decida no campo.
  \item \textbf{Corrija antes de precisar}: Calcário reage em meses; aplicado na janela seca, chega ao plantio já reagido.
  \item \textbf{Chuva em mm não é água no solo}: Textura, declive e cobertura morta decidem quanto do que cai chega à raiz.
  \item \textbf{Série curta é pista, não padrão}: Poucos anos fechados oscilam muito: compare com a estação vizinha antes de fixar calendário.
\end{itemize}

\vspace{2pt}
{\footnotesize\color{cinza}
Fonte: CIIAGRO/IAC (ciiagro.org.br), série mensal medida; chuva é o total do mês.
Painel interativo:
\href{https://climalitoralsul.netlify.app/\#registro-unesp}{\texttt{climalitoralsul.netlify.app/\#registro-unesp}}\\
Umidade relativa ficou fora: a fonte usa \texttt{0} para ``sem sensor'', o que tornaria a média enganosa. Rótulo do painel: médias do período 2009–2026 (15.9 anos de dado; parte da série não validada).
8 leitura(s) fora da faixa de sanidade excluída(s) da média (tmin 3, tmed 3, tmax 2); o banco as preserva.
Gerado a partir de \texttt{dados.json} de 2026-10-30.}

\end{document}
